\FigureLayoutDeclare{img/2004/shanghai_science/q18_fig1.png}{4.0cm}{68bp 15bp 11bp 38bp}

\chapter{2004年上海卷（秋理）}
\section{填空题}
\begin{problem}
若 \(\tan \alpha = \frac{1}{2}\)，则 \(\tan \left(\alpha + \frac{\pi}{4}\right) =\fillinblank{}\).
\end{problem}

\begin{answer}
\(3\)
\end{answer}

\begin{problem}
设抛物线的顶点坐标为 \((2,0)\)，准线方程为 \(x = -1\)，则它的焦点坐标为\fillinblank{}.
\end{problem}

\begin{answer}
\((5,0)\)
\end{answer}

\begin{problem}
设集合 \(A = \{5, \log_2(a + 3)\}\)，集合 \(B = \{a, b\}\). 若 \(A \cap B = \{2\}\)，则 \(A \cup B =\fillinblank{}\).

\end{problem}

\begin{answer}
\(\{1,2,5\}\)
\end{answer}

\begin{problem}
设等比数列 \(\{a_{n}\}\)（\(n \in \N\)）的公比 \(q = -\frac{1}{2}\)，且 \(\displaystyle\lim\limits_{n \to \infty} (a_{1} + a_{3} + a_{5} + \cdots + a_{2n-1}) = \frac{8}{3}\)，则 \(a_{1} =\fillinblank{}\).

\end{problem}

\begin{answer}
\(2\)
\end{answer}

\begin{problem}
设奇函数 \( f(x) \) 的定义域为 \([-5, 5]\). 若当 \( x \in [0, 5] \) 时，\( f(x) \) 的图象如右图，则不等式 \( f(x) < 0 \) 的解是\fillinblank{}.

\bitmapfigure[max width=0.34\linewidth]{img/2004/shanghai_science/q5_fig1.png}
\end{problem}

\begin{answer}
\((-2,0) \cup (2,5]\)
\end{answer}

\begin{problem}
已知点 \(A(1, -2)\)，若向量 \(\overrightarrow{AB}\) 与 \(\bs{a} = (2, 3)\) 同向，\(\left|\overrightarrow{AB}\right| = 2\sqrt{13}\)，则点 \(B\) 的坐标为\fillinblank{}.

\end{problem}

\begin{answer}
\((5,4)\)
\end{answer}

\begin{problem}
在极坐标系中，点 \(M\left(4, \frac{\pi}{3}\right)\) 到直线 \(l: \rho(2\cos\theta + \sin\theta) = 4\) 的距离 \(d =\fillinblank{}\).
\end{problem}

\begin{answer}
\(\frac{2\sqrt{15}}{5}\)
\end{answer}

\begin{problem}
圆心在直线 \(2x - y - 7 = 0\) 上的圆 \(C\) 与 \(y\) 轴交于两点 \(A(0, -4)\)，\(B(0, -2)\)，则圆 \(C\) 的方程为\fillinblank{}.

\end{problem}

\begin{answer}
\(\left(x-2\right)^2+\left(y+3\right)^2=5\)
\end{answer}

\begin{problem}
若在二项式 \((x + 1)^{10}\) 的展开式中任取一项，则该项的系数为奇数的概率是\fillinblank{}.（结果用分数表示）
\end{problem}

\begin{answer}
\(\frac{4}{11}\)
\end{answer}

\begin{problem}
若函数 \( f(x) = a|x - b| + 2 \) 在 \([0, +\infty)\) 上为增函数，则实数 \(a\)，\(b\) 的取值范围是\fillinblank{}.

\end{problem}

\begin{answer}
\(a>0\) 且 \(b\leqslant 0\)
\end{answer}

\begin{problem}
教材中“坐标平面上的直线”与“圆锥曲线”两章内容体现出解析几何的本质是\fillinblank{}.

\end{problem}

\begin{answer}
用代数的方法研究图形的几何性质.
\end{answer}

\begin{problem}
若干个能唯一确定一个数列的量称为该数列的“基本量”. 设 \(\{a_{n}\}\) 是公比为 \(q\) 的无穷等比数列，下列 \(\{a_{n}\}\) 的四组量中，一定能成为该数列“基本量”的是第\fillinblank{} 组. （写出所有符合要求的组号）

\begin{circlelist}
\item \(S_{1}\) 与 \(S_{2}\).
\item \(a_{2}\) 与 \(S_{3}\).
\item \(a_{1}\) 与 \(a_{n}\).
\item \(q\) 与 \(a_{n}\).
\end{circlelist}

其中 \(n\) 为大于 1 的整数，\(S_{n}\) 为 \(\{a_{n}\}\) 的前 \(n\) 项和.
\end{problem}

\begin{answer}
\(\circled{1}\)、\(\circled{4}\)
\end{answer}

\section{单选题}
\begin{problem}
在下列关于直线 \(l\)，\(m\) 与平面 \(\alpha\)，\(\beta\) 的命题中，真命题是（\quad）

\choices
    {若 \(l \subset \beta\) 且 \(\alpha \perp \beta\)，则 \(l \perp \alpha\)}
    {若 \(l \perp \beta\) 且 \(\alpha \parallel \beta\)，则 \(l \perp \alpha\)}
    {若 \(l \perp \beta\) 且 \(\alpha \perp \beta\)，则 \(l \parallel \alpha\)}
    {若 \(\alpha \cap \beta = m\) 且 \(l \parallel m\)，则 \(l \parallel \alpha\)}
\end{problem}

\begin{answer}
B
\end{answer}

\begin{problem}
三角方程 \(2\sin \left(\frac{\pi}{2} - x\right) = 1\) 的解集为（\quad）

\choices
    {\(\left\{x \mid x = 2k\pi + \frac{\pi}{3}, k \in \Z\right\}\)}
    {\(\left\{x \mid x = 2k\pi + \frac{5\pi}{3}, k \in \Z\right\}\)}
    {\(\left\{x \mid x = 2k\pi \pm \frac{\pi}{3}, k \in \Z\right\}\)}
    {\(\{x\mid x = k\pi +(-1)^k,k\in \Z\}\)}
\end{problem}

\begin{answer}
C
\end{answer}

\begin{problem}
若函数 \( y = f(x) \) 的图象可由函数 \( y = \lg (x + 1) \) 的图象绕坐标原点 \(O\) 逆时针旋转 \( \frac{\pi}{2} \) 得到，则 \( f(x) = \)（\quad）

\choices
    {\(10^{-x} - 1\)}
    {\(10^{x} - 1\)}
    {\(1 - 10^{-x}\)}
    {\(1 - 10^{x}\)}
\end{problem}

\begin{answer}
A
\end{answer}

\begin{problem}
某地 2004 年第一季度应聘和招聘人数排行榜前 5 个行业的情况列表如下：

\begin{center}
\begin{tabular}{|c|c|c|c|c|c|}
\hline
行业名称 & 计算机 & 机械 & 营销 & 物流 & 贸易 \\
\hline
应聘人数 & \(215830\) & \(200250\) & \(154676\) & \(74570\) & \(65280\) \\
\hline
\end{tabular}

\medskip

\begin{tabular}{|c|c|c|c|c|c|}
\hline
行业名称 & 计算机 & 营销 & 机械 & 建筑 & 化工 \\
\hline
招聘人数 & \(124620\) & \(102935\) & \(89115\) & \(76516\) & \(70436\) \\
\hline
\end{tabular}
\end{center}

若用同一行业中应聘人数与招聘人数比值的大小来衡量该行业的就业情况，则根据表中数据，就业形势一定是（\quad）

\choices
    {计算机行业好于化工行业}
    {建筑行业好于物流行业}
    {机械行业最紧张}
    {营销行业比贸易行业紧张}
\end{problem}

\begin{answer}
B
\end{answer}

\section{解答题}
\begin{problem}
已知复数 \(z_{1}\) 满足 \((1 + \i)z_{1} = -1 + 5\i\)，\(z_{2} = a - 2 - \i\)，其中 \(\i\) 为虚数单位，\(a \in \R\)，若 \(\left|z_{1} - \overline{z_{2}}\right| < |z_{1}|\)，求 \(a\) 的取值范围.
\end{problem}

\begin{solution}
由 \((1+\i)z_{1}=-1+5\i\)，得
\[
z_{1}=\frac{-1+5\i}{1+\i}=\frac{(-1+5\i)(1-\i)}{2}=2+3\i.
\]
又因为 \(z_{2}=a-2-\i\)，所以 \(\overline{z_{2}}=a-2+\i\)，从而
\[
z_{1}-\overline{z_{2}}=(2+3\i)-(a-2+\i)=4-a+2\i.
\]
因此 \(\left|z_{1}-\overline{z_{2}}\right|^{2}=(4-a)^{2}+4\)，而 \(|z_{1}|^{2}=2^{2}+3^{2}=13\). 题设不等式等价于 \((4-a)^{2}+4<13\)，即 \((a-4)^{2}<9\). 所以 \(1<a<7\).
\end{solution}

\begin{problem}
某单位用木料制作如图所示的框架，框架的下部是边长分别为 \(x\)，\(y\)（单位：\(\mathrm{m}\)）的矩形. 上部是等腰直角三角形. 要求框架围成的总面积 \(8 \mathrm{~cm}^{2}\). 问 \(x\)，\(y\) 分别为多少（精确到 \(0.001 \mathrm{~m}\)）时用料最省？

\bitmapfigure[max width=0.35\linewidth]{img/2004/shanghai_science/q18_fig1.png}
\end{problem}

\begin{solution}
上部等腰直角三角形的底边为斜边，故其两腰长为 \(\dfrac{x}{\sqrt{2}}\)，高为 \(\dfrac{x}{2}\). 因而框架围成的面积条件为 \(xy+\dfrac{x^{2}}{4}=8\)，其中 \(x>0\)，\(y>0\). 于是
\[
y=\frac{8}{x}-\frac{x}{4},\qquad 0<x<4\sqrt{2}.
\]
矩形上边与三角形底边重合，只需计算框架的外边木料长度，所以用料长度为
\[
L=2x+2y+2\cdot\frac{x}{\sqrt{2}}=\left(\frac{3}{2}+\sqrt{2}\right)x+\frac{16}{x}.
\]
由基本不等式，
\[
L\geqslant2\sqrt{16\left(\frac{3}{2}+\sqrt{2}\right)},
\]
当且仅当 \(\left(\dfrac{3}{2}+\sqrt{2}\right)x=\dfrac{16}{x}\) 时取等号. 因此
\[
x=\frac{4}{\sqrt{\frac{3}{2}+\sqrt{2}}}=8-4\sqrt{2},
\qquad y=\frac{8}{x}-\frac{x}{4}=2\sqrt{2}.
\]
所以当 \(x\approx2.343\mathrm{~m}\)，\(y\approx2.828\mathrm{~m}\) 时用料最省.
\end{solution}

\begin{problem}
记函数 \( f(x) = \sqrt{2 - \frac{x + 3}{x + 1}} \) 的定义域为 \(A\)，\( g(x) = \lg [(x - a - 1)(2a - x)]\)（\(a < 1\)）的定义域为 \(D\).

\begin{enumerate}[itemsep=0pt,topsep=0pt]
\item 求 \(A\).
\item 若 \(B \subseteq A\)，求实数 \(a\) 的取值范围.
\end{enumerate}
\end{problem}

\begin{solution}
\begin{enumerate}
\item 由根式有意义，且 \(x\neq-1\)，得
\[
2-\frac{x+3}{x+1}=\frac{x-1}{x+1}\geqslant0.
\]
所以 \(x<-1\) 或 \(x\geqslant1\)，即 \(A=(-\infty,-1)\cup[1,+\infty)\).

\item 题干把 \(g(x)\) 的定义域记为 \(D\)，第（2）问却写作 \(B\subseteq A\). 以下按题意将第（2）问中的 \(B\) 解释为该定义域 \(D\). 由对数真数为正，得
\[
(x-a-1)(2a-x)>0.
\]
因为 \(a<1\)，所以 \(2a<a+1\)，从而 \(D=(2a,a+1)\). 要使 \(D\subseteq(-\infty,-1)\cup[1,+\infty)\)，开区间 \(D\) 必须完全落在左侧区间或右侧区间，因此
\[
a+1\leqslant-1\quad\text{或}\quad 2a\geqslant1.
\]
结合 \(a<1\)，得到 \(a\leqslant-2\) 或 \(\dfrac{1}{2}\leqslant a<1\). 故实数 \(a\) 的取值范围为 \((-\infty,-2]\cup[\dfrac{1}{2},1)\).
\end{enumerate}
\end{solution}

\begin{problem}
已知二次函数 \( y = f_{1}(x) \) 的图象以原点为顶点且过点 \((1,1)\). 反比例函数 \( y = f_{2}(x) \) 的图象与直线 \( y = x \) 的两个交点间距离为 \(8\)，\( f(x) = f_{1}(x) + f_{2}(x) \).

\begin{enumerate}[itemsep=0pt,topsep=0pt]
\item 求函数 \( f(x) \) 的表达式.
\item 证明：当 \(a > 3\) 时，关于 \(x\) 的方程 \(f(x) = f(a)\) 有三个实数解.
\end{enumerate}
\end{problem}

\begin{solution}
\begin{enumerate}
\item 二次函数图象的顶点为原点，故可设 \(f_{1}(x)=cx^{2}\). 又因图象过 \((1,1)\)，有 \(c=1\)，所以 \(f_{1}(x)=x^{2}\).

设反比例函数为 \(f_{2}(x)=\dfrac{k}{x}\)，其中 \(k\neq0\). 它与直线 \(y=x\) 的交点满足 \(x=\dfrac{k}{x}\)，即 \(x^{2}=k\). 题设有两个交点，所以 \(k>0\)，两交点为 \((\sqrt{k},\sqrt{k})\) 和 \((-\sqrt{k},-\sqrt{k})\). 由两点间距离为 \(8\)，得 \(2\sqrt{2k}=8\)，所以 \(k=8\). 因而
\[
f(x)=x^{2}+\frac{8}{x},\qquad x\neq0.
\]

\item 当 \(a>3\) 时，\(a\neq0\). 方程 \(f(x)=f(a)\) 中也必须有 \(x\neq0\)，且
\[
0=f(x)-f(a)=x^{2}-a^{2}+\frac{8}{x}-\frac{8}{a}=(x-a)\left(x+a-\frac{8}{ax}\right).
\]
所以一个实数解为 \(x=a\)，其余解满足
\[
ax^{2}+a^{2}x-8=0.
\]
这个二次方程的判别式为 \(a^{4}+32a>0\)，且两根之积为 \(-\dfrac{8}{a}<0\)，故它有一个负根和一个正根. 正根不可能等于 \(a\)，因为若正根为 \(a\)，则 \(2a=\dfrac{8}{a^{2}}\)，即 \(a^{3}=4\)，这与 \(a>3\) 矛盾. 因此 \(x=a\) 与二次方程的两个根互不相同，方程 \(f(x)=f(a)\) 有三个实数解.
\end{enumerate}
\end{solution}

\begin{problem}
如图，\(P - ABC\) 是底面边长为 \(1\) 的正三棱锥，\(D\)，\(E\)，\(F\) 分别为棱长 \(PA\)，\(PB\)，\(PC\) 上的点，截面 \(DEF \parallel\) 底面 \(ABC\)，且棱台 \(DEF - ABC\) 与棱锥 \(P - ABC\) 的棱长和相等.（棱长和是指多面体中所有棱的长度之和）

\begin{enumerate}[itemsep=0pt,topsep=0pt]
\item 证明：\(P - ABC\) 为正四面体.
\item 若 \(PD = \frac{1}{2} PA\)，求二面角 \(D - BC - A\) 的大小.（结果用反三角函数值表示）
\item 设棱台 \(DEF - ABC\) 的体积为 \(V\)，是否存在体积为 \(V\) 且各棱长均相等的直平行六面体，使得它与棱台 \(DEF - ABC\) 有相同的棱长和？若存在，请具体构造出这样的一个直平行六面体，并给出证明；若不存在，请说明理由.
\end{enumerate}

\bitmapfigure[max width=0.42\linewidth]{img/2004/shanghai_science/q21_fig1.png}
\end{problem}

\begin{solution}
设 \(\dfrac{PD}{PA}=\dfrac{PE}{PB}=\dfrac{PF}{PC}=t\)，由截面平行于底面且位于棱锥内部，得 \(0<t<1\). 设正三棱锥的侧棱长为 \(h\). 由相似三角形，\(DEF\) 为边长为 \(t\) 的正三角形，且 \(AD=BE=CF=(1-t)h\).

\begin{enumerate}
\item 棱锥 \(P-ABC\) 的棱长和为 \(3+3h\)，棱台 \(DEF-ABC\) 的棱长和为
\[
3+3t+3(1-t)h.
\]
由题设两者相等，得 \(3+3t+3(1-t)h=3+3h\)，即 \(t(1-h)=0\). 因为 \(t>0\)，所以 \(h=1\). 底面 \(ABC\) 是边长为 \(1\) 的正三角形，且 \(PA=PB=PC=1\)，故 \(P-ABC\) 的六条棱均相等，是正四面体.

\item 取 \(M\) 为 \(BC\) 的中点，连接 \(PM\)，\(AM\)，\(DM\). 在正三角形 \(ABC\) 和正三角形 \(PBC\) 中，均有 \(BC\perp AM\) 且 \(BC\perp PM\)，所以 \(BC\perp\) 平面 \(PAM\)，从而 \(BC\perp DM\). 因此二面角 \(D-BC-A\) 的平面角为 \(\angle DMA\).

当 \(PD=\dfrac12PA\) 时，\(D\) 是 \(PA\) 的中点. 又 \(AM=PM=\dfrac{\sqrt{3}}{2}\)，\(AP=1\)，由三角形中线公式得 \(DM^{2}=\dfrac12\)，而 \(DA=\dfrac12\)，\(AM^{2}=\dfrac34\). 在三角形 \(DMA\) 中，
\[
\cos\angle DMA=\frac{DM^{2}+AM^{2}-DA^{2}}{2DM\cdot AM}=\sqrt{\frac{2}{3}}.
\]
所以二面角 \(D-BC-A\) 的大小为 \(\arccos\sqrt{\dfrac{2}{3}}=\arcsin\dfrac{1}{\sqrt{3}}\).

\item 由第（2）问，\(t=\dfrac12\). 因而棱台的棱长和为
\[
3+3\cdot\frac12+3\cdot\frac12=6.
\]
边长为 \(1\) 的正四面体体积为 \(\dfrac{\sqrt{2}}{12}\)，而棱锥 \(P-DEF\) 与棱锥 \(P-ABC\) 的相似比为 \(\dfrac12\)，所以
\[
V=\left(1-\left(\frac12\right)^{3}\right)\frac{\sqrt{2}}{12}=\frac{7\sqrt{2}}{96}.
\]

构造一个直平行六面体：取底面为边长 \(\dfrac12\)、一个内角为 \(\alpha=\arcsin\dfrac{7\sqrt{2}}{12}\) 的菱形，取垂直于底面的侧棱，且所有侧棱长均为 \(\dfrac12\). 这是一个各棱长均为 \(\dfrac12\) 的直平行六面体，其棱长和为 \(12\cdot\dfrac12=6\)，体积为
\[
\left(\frac12\right)^{3}\sin\alpha=\frac18\cdot\frac{7\sqrt{2}}{12}=\frac{7\sqrt{2}}{96}=V.
\]
因此这样的直平行六面体存在，以上构造满足全部条件.
\end{enumerate}
\end{solution}

\begin{problem}
设 \(P_{1}(x_{1},y_{1})\)，\(P_{2}(x_{2},y_{2})\)，\(\cdots\)，\(P_{n}(x_{n},y_{n})\)（\(n \geqslant 3\)，\(n \in \N\)）是二次曲线 \(C\) 上的点，且 \(a_{1} = |OP_{1}|^{2}\)，\(a_{2} = |OP_{2}|^{2}\)，\(\cdots\)，\(a_{n} = |OP_{n}|^{2}\) 构成了一个公差为 \(d\)（\(d \neq 0\)）的等差数列，其中 \(O\) 是坐标原点. 记 \(S_{n} = a_{1} + a_{2} + \cdots + a_{n}\).

\begin{enumerate}[itemsep=0pt,topsep=0pt]
\item 若 \(C\) 的方程为 \(\frac{x^2}{100} + \frac{y^2}{25} = 1\)，\(n = 3\). 点 \(P_1(3,0)\) 及 \(S_3 = 255\)，求点 \(P_3\) 的坐标.（只需写出一个）
\item 若 \(C\) 的方程为 \(\frac{x^2}{a^2} + \frac{y^2}{b^2} = 1\)（\(a > b > 0\)）. 点 \(P_1(a,0)\)，对于给定的自然数 \(n\)，当公差 \(d\) 变化时，求 \(S_n\) 的最小值.

\item 请选定一条除椭圆外的二次曲线 \(C\) 及 \(C\) 上的一点 \(P_1\)，对于给定的自然数 \(n\)，写出符合条件的点 \(P_1\)，\(P_2\)，\(\ldots\)，\(P_n\) 存在的充要条件，并说明理由.
\end{enumerate}
\end{problem}

\begin{solution}
\begin{enumerate}
\item 由 \(P_{1}=(3,0)\)，得 \(a_{1}=|OP_{1}|^{2}=3^{2}=9\). 另外，\((3,0)\) 不满足 \(\dfrac{x^{2}}{100}+\dfrac{y^{2}}{25}=1\)，题面前提本身已经不一致. 即使暂不考虑这一点，三项等差数列的中项为平均值，所以
\[
a_{2}=\frac{S_{3}}{3}=\frac{255}{3}=85,
\qquad a_{3}=2a_{2}-a_{1}=2\cdot85-9=161.
\]
但若 \(P_{3}=(x_{3},y_{3})\) 在椭圆 \(\dfrac{x^{2}}{100}+\dfrac{y^{2}}{25}=1\) 上，则
\[
x_{3}^{2}+4y_{3}^{2}=100,
\qquad |OP_{3}|^{2}=x_{3}^{2}+y_{3}^{2}=100-3y_{3}^{2}\leqslant100.
\]
这与 \(a_{3}=|OP_{3}|^{2}=161\) 矛盾. 因此按原试卷印出的数据，不存在满足条件的点 \(P_{3}\).

\item 在椭圆 \(\dfrac{x^{2}}{a^{2}}+\dfrac{y^{2}}{b^{2}}=1\) 上，有
\[
x^{2}+y^{2}=b^{2}+\left(1-\frac{b^{2}}{a^{2}}\right)x^{2},
\]
所以原点到椭圆上点的距离平方的取值范围为 \([b^{2},a^{2}]\). 由于 \(P_{1}=(a,0)\)，有 \(a_{1}=a^{2}\)，它已经是最大值. 又因 \(d\neq0\)，要使 \(a_{2}=a^{2}+d\leqslant a^{2}\)，必须有 \(d<0\).

数列末项也必须满足 \(a_{n}=a^{2}+(n-1)d\geqslant b^{2}\)，故
\[
\frac{b^{2}-a^{2}}{n-1}\leqslant d<0.
\]
反之，任取该区间内的 \(d\)，各项 \(a_{i}=a^{2}+(i-1)d\) 均属于 \([b^{2},a^{2}]\)，都可在该椭圆上取到，因此这正是允许的公差范围. 于是
\[
S_{n}=na^{2}+\frac{n(n-1)}{2}d.
\]
由于 \(\frac{n(n-1)}{2}>0\)，\(S_{n}\) 随 \(d\) 增大而增大，故在 \(d=\frac{b^{2}-a^{2}}{n-1}\) 时取得最小值：
\[
S_{n,\min}=na^{2}+\frac{n(n-1)}{2}\cdot\frac{b^{2}-a^{2}}{n-1}=\frac{n(a^{2}+b^{2})}{2}.
\]

\item 取双曲线 \(C\colon x^{2}-y^{2}=1\)，并取 \(P_{1}=(1,0)\). 对任意 \(P=(x,y)\in C\)，有
\[
|OP|^{2}=x^{2}+y^{2}=1+2y^{2}\geqslant1.
\]
而 \(a_{1}=|OP_{1}|^{2}=1\) 已是最小值. 因此若 \(P_{1},P_{2},\ldots,P_{n}\) 存在，由 \(d\neq0\) 及 \(a_{2}=1+d\geqslant1\) 必有 \(d>0\).

反之，若 \(d>0\)，令 \(a_{i}=1+(i-1)d\)，并取
\[
P_{i}=\left(\sqrt{\frac{a_{i}+1}{2}},\sqrt{\frac{a_{i}-1}{2}}\right)\quad(i=1,2,\ldots,n).
\]
则 \(x_{i}^{2}-y_{i}^{2}=1\)，所以 \(P_{i}\in C\)，且 \(|OP_{i}|^{2}=a_{i}\). 因此满足条件的点存在的充要条件为 \(d>0\).
\end{enumerate}
\end{solution}
